% =====================================================================
% Sprint 2, Track 7: Gauge-choice projection (§III.26)
%
% Placement: §III.26, immediately after §III.25 coupled-channel /
% adiabatic curve (label sec:proj_coupled_channel), immediately before
% the §III boundary subsection (label sec:layer2_boundary).
%
% §IV three-axis row (for tab:projection_axes, insert after the
% coupled-channel row, last row before \bottomrule):
% ---------------------------------------------------------------------
%   Gauge choice (Coulomb / Lorenz / Feynman) &
%     none (gauge choice, not parameter) &
%     preserves &
%     gauge-dependent at intermediate-propagator level
%       (Coulomb: $1/(4\pi)$ per loop via Hopf base;
%        Lorenz: retarded; Feynman: $1/k^2$);
%       gauge-invariant on physical observables (Ward identity) \\
% ---------------------------------------------------------------------
% One-sentence summary for §IV / count update:
%   "The twenty-sixth row (gauge choice) was added in Sprint 2
%    (2026-05-15) as the second witness, alongside vector-photon
%    promotion, in the (no-variable, no-dimension, transcendental-only)
%    corner of the V/D/transcendental axis grid."
%
% Cross-references inside this draft:
%   - sec:proj_vector_photon (§III.11): first independence witness
%   - sec:proj_3j (§III.8): Wigner-3j angular coupling
%   - sec:proj_wilson (§III.10): Wilson plaquette (implicit Coulomb-class
%     gauge fixing on the GeoVac SU(2)/U(1) Wilson construction)
%   - sec:proj_spectral_action (§III.6): heat-kernel / spectral action
%   - sec:proj_multipole_gaunt (§III.21): Coulomb kernel multipole
%     expansion (angular-content-only, same V/D-preserving family)
%   - sec:layer2 (the projection-family section)
%   - sec:layer2_boundary (the boundary subsection that follows)
%
% Cross-references to other papers (use existing citation keys):
%   - Paper~28 (one-loop QED machinery, Coulomb-gauge construction)
%   - Paper~33 (1/(4\pi) per loop = S^2 Weyl exchange of Hopf base)
%   - Paper~30 (SU(2) Wilson lattice; gauge fixing implicit in
%     weak-coupling kinetic identification L_1 = B^\top B)
%   - Paper~18 (calibration tier, in which the per-loop 1/(4\pi)
%     lives)
%
% Observation update needed in §IV (rem:vd_correlation /
% obs:axes_correlated): the transcendental-only corner now has
% two witnesses (vector-photon promotion + gauge choice), not one.
% Suggested wording: "The transcendental axis admits two
% independence witnesses --- vector-photon promotion and gauge
% choice (\S~\ref{sec:proj_gauge_choice}), both of which preserve
% variable and dimension but inject transcendental content (or
% structurally select which transcendental content appears) at
% intermediate steps.  Both witnesses are gauge-related: the
% structural-content of the transcendental-only corner is the
% gauge sector of the framework."
% =====================================================================

\subsection{Gauge choice (Coulomb / Lorenz / Feynman)}
\label{sec:proj_gauge_choice}
\textbf{Source:} the framework's $U(1)$ connection on the Fock-projected
$S^3$ Hopf graph (with $\mathrm{SU}(2)$ and $\mathrm{SU}(3)$ siblings
on the same graph and on $S^5$ respectively, per Paper~30 and
Sprint~ST-SU3); equivalently, the photon $1$-cochain in the Hodge
decomposition of edge data, prior to any specific gauge representative
being selected for propagator computations.\\
\textbf{Target:} the same connection with a specific gauge
representative fixed.  Three canonical choices are in play:
\begin{itemize}\setlength{\itemsep}{2pt}
\item \emph{Coulomb gauge} ($\vec{\nabla} \cdot \vec{A} = 0$,
equivalently $d^\dagger A = 0$ on the graph): the photon decomposes
into an instantaneous longitudinal Coulomb part and a transverse
radiative part; the framework's vector-photon construction
(Paper~28~\S vector\_photon\_qed, Paper~33~\S~VI) operates in this
gauge by treating the Hopf-base $1$-form sector as the propagating
photon and computing $G_\gamma = L_1^+$ as the transverse propagator
(see also \S~\ref{sec:proj_wilson}, where the weak-coupling expansion
$L_1 = B^\top B$ is a Coulomb-class gauge fixing of the SU(2) Wilson
construction).
\item \emph{Lorenz gauge} ($\partial_\mu A^\mu = 0$, equivalently the
covariant retarded condition on the spacetime $4$-potential): the
photon propagator carries explicit retardation in time and the
transverse and longitudinal pieces remain coupled.
\item \emph{Feynman gauge} (the propagator $-i g_{\mu\nu}/k^2$): the
simplest covariant choice, used in Paper~28~\S~6 for the one-loop
electron self-energy (``the one-loop electron self-energy on $S^3$
(Feynman gauge) is \ldots'').
\end{itemize}
The projection step is the act of fixing one of these representatives
\emph{before} intermediate Green's-function, propagator, or matrix-element
computations begin.\\
\textbf{Variables introduced:} \emph{none}.  A gauge choice is a
selection among equivalent classes of representatives, not the
introduction of a continuous parameter.  No physical scale,
focal-length, mass, or coordinate is injected by the gauge fixing
itself.  This places gauge choice alongside \S~\ref{sec:proj_vector_photon}
in the (no-variable) row of Table~\ref{tab:projection_axes}.\\
\textbf{Dimension introduced:} dimension-preserving.  The gauge choice
operates on the existing connection without injecting a new $[L]$,
$[E]$, $[M]$, or $[T]$ scale.  All three gauges produce the same
physical observable with the same physical dimension; what differs is
the intermediate-step organisation of the propagator content.\\
\textbf{Transcendental signature:} \emph{gauge-dependent at
intermediate steps; gauge-invariant on physical observables.}  This is
the distinguishing structural property of the gauge-choice projection.
At the intermediate-propagator level, the three gauges produce
categorically different transcendental content:
\begin{itemize}\setlength{\itemsep}{2pt}
\item In Coulomb gauge --- the framework's implicit choice throughout
Papers~28, 30, and 33 --- the vector-photon construction's per-loop
factor is $1/(4\pi) = \mathrm{Vol}(S^2)/(4 \cdot 4\pi^2)$, identified
in Paper~33~\S~VI as the $S^2$ Weyl exchange constant of the Hopf
base (the same Hopf-base measure $\pi = \mathrm{Vol}(S^2)/4$ that
appears in the M1 master-Mellin sub-mechanism, per Paper~32~\S~VIII
and Paper~18~\S~III.7).  The intermediate transcendental content sits
cleanly in the calibration $\pi$-tier of Paper~18.
\item In Lorenz gauge, the retarded propagator structure would produce
different intermediate angular content, with $\pi$ entering through
the retardation kernel rather than through the static Hopf-base solid
angle.  This route has not been implemented in the framework.
\item In Feynman gauge, the explicit $-i g_{\mu\nu}/k^2$ propagator
gives the textbook covariant intermediate transcendental content, with
$\pi$ appearing through the standard Feynman-parameter denominators
rather than through a Hopf-base normalisation.  Paper~28~\S~6 names
this gauge for the one-loop self-energy formula but the production
machinery operates in Coulomb gauge throughout.
\end{itemize}
On physical observables (cross-sections, transition rates, energy
shifts), all three gauges produce identical results --- this is the
content of the Ward identity (Rule~6 of Paper~33's eight QED selection
rules; see Paper~33~\S~II for the framework's status), and it is
gauge-invariance in its standard QED sense.\\
\textbf{Used in:} the framework's one-loop QED machinery (Paper~28)
implicitly assumes Coulomb gauge throughout the $1$-cochain Hodge
construction; Paper~33's $1{+}6{+}1$ selection-rule partition is
stated in Coulomb gauge; Paper~30's SU(2) Wilson lattice gauge theory
identifies $L_1 = B^\top B$ as the weak-coupling kinetic term via a
Coulomb-class gauge fixing of the link variables.  Feynman gauge is
named at one point in Paper~28~\S~6 (the textbook formula for the
one-loop electron self-energy) but is not used in production
calculations.  No production computation has been performed in Lorenz
gauge.\\
\textbf{Empirical anchor:} the Ward identity itself, i.e.\
gauge-invariance of physical observables.  Standard QED requires
$k_\mu \mathcal{M}^\mu = 0$ for the photon-emission amplitude
$\mathcal{M}^\mu$, and equivalently the equality of physical
predictions across gauge choices.  This is item~R6 in Paper~33's
selection-rule census, where the framework's status is recorded
honestly: ``gauge invariance under the $U(1)$ phase of the photon,
which the $1$-cochain photon has only in temporal gauge and only
approximately'' (Paper~33~\S~II, R6).  The Ward identity is therefore
the standard external anchor for the gauge-choice projection;
gauge-invariance of physical observables is the precise condition
under which the projection itself is structurally well-posed.\\
\textbf{Honest scope.}  This is the single most under-tested entry in
the projection dictionary, and the entry must be tagged as such.

The framework has \emph{not} been tested across gauges.  Every
production calculation in Papers~28, 30, 33, and the Sprint~HF /
Sprint~LS-1..LS-7 / Lamb-shift sub-percent closure of Paper~36 runs
implicitly in Coulomb gauge: the photon is represented as a Hodge
$1$-cochain on the GeoVac edge complex, the propagator is
$G_\gamma = L_1^+$ (the transverse pseudoinverse of the edge
Laplacian), and the per-loop calibration constant $1/(4\pi)$ is
identified through the Hopf-base $S^2$ solid-angle normalisation.
This is a clean, structurally well-motivated choice ---  Coulomb gauge
respects the framework's natural angular-momentum decomposition on
the Fock-projected $S^3$ graph, and the transverse photon sector
sits naturally in the co-exact subspace of the Hodge decomposition
of \S~\ref{sec:proj_wilson} (compare Paper~28~\S~transverse\_photon).
But it \emph{is} a choice.

A genuine validation of the gauge-choice projection at framework level
would require computing a sample one-loop observable in Lorenz gauge
(or Feynman gauge as implemented at full operator level rather than
named for one formula), verifying that the physical result agrees
with the Coulomb-gauge calculation to the precision of the framework,
and confirming that the per-loop calibration constant changes
structurally between the two gauges while the physical observable is
gauge-invariant.  This has not been done.  The Ward identity (R6 in
Paper~33's census) holds only approximately on the $1$-cochain photon
even in temporal gauge, and the framework has not closed the
gauge-invariance test against a competing gauge.

Sub-leading observation: the choice of Coulomb gauge is not arbitrary
on the GeoVac graph.  The Hodge decomposition of the edge complex
distinguishes exact ($d_0$-image), co-exact ($d_1^\top$-image), and
harmonic sectors, and Coulomb gauge is precisely the projection onto
the co-exact (transverse) sector.  The framework's natural data
structure --- $L_0$ on nodes, $L_1 = B^\top B$ on edges,
plaquette Hodge structure from Paper~25 / Paper~30 --- comes equipped
with this decomposition.  Lorenz gauge, by contrast, would require
introducing the time direction explicitly (compactified per
\S~\ref{sec:proj_observation} or extended per the Klein-Gordon
construction of Paper~35), and the framework has not been wired for
covariant time-dependent gauge fixing.  This is a real structural
constraint, not a choice that could be flipped without modifying the
upstream construction.\\
\textbf{Structural reading.}  Layer~$2 \to$ Layer~$2$.  The gauge-choice
projection operates on the projected vector-photon Hilbert space
(\S~\ref{sec:proj_vector_photon}); it does not act on Layer~1 graph
data and does not move the chain across the
\S~\ref{sec:layer2_boundary} boundary.

This is the \emph{second independence witness} in the
(no-variable, no-dimension, transcendental-only) corner of the
variable/dimension/transcendental axis grid.  Until Sprint~2 the only
witness in that corner was \S~\ref{sec:proj_vector_photon} (vector-photon
promotion, which introduces $1/(4\pi)$ without any new variable or
dimension); the gauge-choice projection is its structural sibling.
Both witnesses operate on the same object (the framework's $U(1)$
connection) at successive structural steps: vector-photon promotion
transforms the scalar $1$-cochain into a vector harmonic with
$(q, m_q)$ labels, and the present projection fixes a specific gauge
representative of the resulting connection before propagators are
computed.  The fact that the two witnesses are \emph{both
gauge-related} sharpens the V/D-correlation observation of
\S~\ref{sec:dictionary} and \S~\ref{sec:axiomatization}: the
transcendental-only corner of the axis grid is the gauge sector of the
framework.  It is not an accident that the only two projections
preserving both V and D while moving the transcendental axis are
gauge-construction steps; the structural content of that corner is
``add transcendental content to the gauge sector without adding any
physical variable or dimension.''

Categorically distinct from \S~\ref{sec:proj_vector_photon} despite
the shared axis pattern: vector-photon promotion is the structural
\emph{construction} of the gauge sector (it builds the vector-harmonic
photon from a scalar $1$-cochain), while the gauge-choice projection
is the \emph{selection} of a representative within the already-built
sector.  The two compose: the chain $\textsf{Layer~1 graph} \to
\S~\ref{sec:proj_vector_photon} \to \S~\ref{sec:proj_gauge_choice}$ is
how the framework reaches a computable QED propagator from bare graph
data.  Categorically distinct also from
\S~\ref{sec:proj_multipole_gaunt}, which similarly preserves V and D
(the multipole index $L$ is an internal integer summation label, not
a physical variable; the kernel reorganisation is dimension-preserving)
but lives in the kernel-side angular sector rather than the gauge
sector and is ring-preserving over $\mathbb{Q}[\sqrt{2k{+}1}]_k$ rather
than gauge-dependent in transcendental content.  The multipole / Gaunt
projection is an angular-content-only structural projection on the
Coulomb kernel; the present projection is a gauge-content-only
structural projection on the connection that the Coulomb kernel
contracts against.  Both are necessary to reach a computable
amplitude; both preserve the V/D axes; neither introduces a Layer-2
input.

Sibling in the calibration tier of Paper~18: the per-loop $1/(4\pi)$
that appears in Coulomb gauge is the same calibration $\pi$-content
identified by Paper~33~\S~VI as the $S^2$ Weyl exchange constant.
Under the master-Mellin reading of Paper~32~\S~VIII and
Paper~18~\S~III.7, this calibration $\pi$ is the M1 sub-mechanism
signature ($\mathrm{Vol}(S^2)/4$, the Hopf-base measure).  A
Lorenz-gauge calculation, if performed, would produce different
intermediate $\pi$-content (retardation kernels rather than static
Hopf-base normalisations) but would necessarily land on the same
physical observable.  The gauge-choice projection therefore sits at
the interface of the framework's gauge construction
(\S~\ref{sec:proj_vector_photon} +
\S~\ref{sec:proj_wilson}) and the framework's transcendental taxonomy
(Paper~18), and it is the structural step at which a Layer-2 chain
commits to a specific way of accumulating $\pi$ at each loop.

Open question (flagged for a future sprint): a serious test of the
gauge-choice projection would compute a sample observable (e.g.\ the
one-loop Lamb-shift Bethe-log component of Paper~36, or the
hydrogen $21$\,cm contact term of \S~\ref{sec:autopsy_21cm}) in
Feynman gauge as well as the existing Coulomb-gauge implementation,
and verify gauge invariance to the precision of the framework
($\sim10^{-3}$ on Lamb-shift-class observables, $\sim10^{-4}$ on
hyperfine contact terms).  If gauge invariance fails at framework
precision, the failure localises a specific Layer-2 calibration
content that the Coulomb-gauge construction is silently encoding;
if it succeeds, the gauge-choice projection is verified as a
genuine independence witness in the sense of
Observation~\ref{obs:axes_correlated} rather than a tautological
relabelling.
